\documentclass[10pt,a4paper]{amsart}
\usepackage{amsmath,amssymb,amsthm,graphicx}
\usepackage{verbatim}
\newenvironment{reference}{\comment}{\endcomment}
%\newenvironment{reference}{}{}
\newenvironment{slogan}{\comment}{\endcomment}
\newenvironment{history}{\comment}{\endcomment}
% For commutative diagrams we use Xy-pic
\usepackage[all]{xy}

% We use 2cell for 2-commutative diagrams.
\xyoption{2cell}
% Theorem environments.
%
\theoremstyle{plain}
\newtheorem{theorem}[subsection]{Theorem}
\newtheorem{proposition}[subsection]{Proposition}
\newtheorem{lemma}[subsection]{Lemma}

\theoremstyle{definition}
\newtheorem{definition}[subsection]{Definition}
\newtheorem{example}[subsection]{Example}
\newtheorem{exercise}[subsection]{Exercise}
\newtheorem{situation}[subsection]{Situation}

\theoremstyle{remark}
\newtheorem{remark}[subsection]{Remark}
\newtheorem{remarks}[subsection]{Remarks}

\numberwithin{equation}{subsection}

% Macros
%
\def\lim{\mathop{\mathrm{lim}}\nolimits}
\def\colim{\mathop{\mathrm{colim}}\nolimits}
\def\Spec{\mathop{\mathrm{Spec}}}
\def\Hom{\mathop{\mathrm{Hom}}\nolimits}
\def\Ext{\mathop{\mathrm{Ext}}\nolimits}
\def\SheafHom{\mathop{\mathcal{H}\!\mathit{om}}\nolimits}
\def\SheafExt{\mathop{\mathcal{E}\!\mathit{xt}}\nolimits}
\def\Sch{\mathit{Sch}}
\def\Mor{\mathop{\mathrm{Mor}}\nolimits}
\def\Ob{\mathop{\mathrm{Ob}}\nolimits}
\def\Sh{\mathop{\mathit{Sh}}\nolimits}
\def\NL{\mathop{N\!L}\nolimits}
\def\CH{\mathop{\mathrm{CH}}\nolimits}
\def\proetale{{pro\text{-}\acute{e}tale}}
\def\etale{{\acute{e}tale}}
\def\QCoh{\mathit{QCoh}}
\def\Ker{\mathop{\mathrm{Ker}}}
\def\Im{\mathop{\mathrm{Im}}}
\def\Coker{\mathop{\mathrm{Coker}}}
\def\Coim{\mathop{\mathrm{Coim}}}

% Boxtimes
%
\DeclareMathSymbol{\boxtimes}{\mathbin}{AMSa}{"02}

%
% Macros for moduli stacks/spaces
%
\def\QCohstack{\mathcal{QC}\!\mathit{oh}}
\def\Cohstack{\mathcal{C}\!\mathit{oh}}
\def\Spacesstack{\mathcal{S}\!\mathit{paces}}
\def\Quotfunctor{\mathrm{Quot}}
\def\Hilbfunctor{\mathrm{Hilb}}
\def\Curvesstack{\mathcal{C}\!\mathit{urves}}
\def\Polarizedstack{\mathcal{P}\!\mathit{olarized}}
\def\Complexesstack{\mathcal{C}\!\mathit{omplexes}}
% \Pic is the operator that assigns to X its picard group, usage \Pic(X)
% \Picardstack_{X/B} denotes the Picard stack of X over B
% \Picardfunctor_{X/B} denotes the Picard functor of X over B
\def\Pic{\mathop{\mathrm{Pic}}\nolimits}
\def\Picardstack{\mathcal{P}\!\mathit{ic}}
\def\Picardfunctor{\mathrm{Pic}}
\def\Deformationcategory{\mathcal{D}\!\mathit{ef}}
\usepackage{cleveref}
\crefname{originaltheorem}{original theorem}{original theorems}
\Crefname{originaltheorem}{Original theorem}{Original theorems}
\makeatletter
\newlabel{algebra-theorem-lazard}{{058G}{source}{}{sourcepointer-1}{}}
\newlabel{algebra-theorem-lazard@cref}{{[originaltheorem][0][]058G}{[1][0][]source}{}{}{}}
\makeatother
\crefname{originallemma}{original lemma}{original lemmas}
\Crefname{originallemma}{Original lemma}{Original lemmas}
\makeatletter
\newlabel{lemma-affine-vanishing-higher-direct-images}{{073H}{source}{}{sourcepointer-2}{}}
\newlabel{lemma-affine-vanishing-higher-direct-images@cref}{{[originallemma][0][]073H}{[1][0][]source}{}{}{}}
\makeatother
\crefname{originallemma}{original lemma}{original lemmas}
\Crefname{originallemma}{Original lemma}{Original lemmas}
\makeatletter
\newlabel{lemma-colimits}{{073E}{source}{}{sourcepointer-3}{}}
\newlabel{lemma-colimits@cref}{{[originallemma][0][]073E}{[1][0][]source}{}{}{}}
\makeatother
\crefname{originallemma}{original lemma}{original lemmas}
\Crefname{originallemma}{Original lemma}{Original lemmas}
\makeatletter
\newlabel{sites-cohomology-lemma-apply-Leray}{{0733}{source}{}{sourcepointer-4}{}}
\newlabel{sites-cohomology-lemma-apply-Leray@cref}{{[originallemma][0][]0733}{[1][0][]source}{}{}{}}
\makeatother
\crefname{originallemma}{original lemma}{original lemmas}
\Crefname{originallemma}{Original lemma}{Original lemmas}
\makeatletter
\newlabel{sites-cohomology-lemma-base-change-map-flat-case}{{0736}{source}{}{sourcepointer-5}{}}
\newlabel{sites-cohomology-lemma-base-change-map-flat-case@cref}{{[originallemma][0][]0736}{[1][0][]source}{}{}{}}
\makeatother
\crefname{originallemma}{original lemma}{original lemmas}
\Crefname{originallemma}{Original lemma}{Original lemmas}
\makeatletter
\newlabel{spaces-morphisms-lemma-base-change-flat}{{03MO}{source}{}{sourcepointer-6}{}}
\newlabel{spaces-morphisms-lemma-base-change-flat@cref}{{[originallemma][0][]03MO}{[1][0][]source}{}{}{}}
\makeatother
\crefname{originallemma}{original lemma}{original lemmas}
\Crefname{originallemma}{Original lemma}{Original lemmas}
\makeatletter
\newlabel{spaces-morphisms-lemma-flat-morphism-sites}{{073C}{source}{}{sourcepointer-7}{}}
\newlabel{spaces-morphisms-lemma-flat-morphism-sites@cref}{{[originallemma][0][]073C}{[1][0][]source}{}{}{}}
\makeatother
\title{Flat base change for the cohomology of a quasi-coherent module on an algebraic space}
\author{The Stacks Project Authors}
\date{stacks-project:spaces-cohomology.tex:local-snapshot-20261012; GNU FDL 1.2 or later}
\begin{document}
\maketitle
\noindent\textit{Editorial note.} All delivered mathematical text is the author's own, copied from the pinned Stacks Project chapter file spaces-cohomology.tex; no statement or proof step was rewritten, reordered or removed.\par\medskip
\noindent\textit{Editorial note.} Preamble provenance: the pinned chapter file opens with \textbackslash{}input\{preamble\} and carries no preamble text of its own, so the project's single Stacks-wide preamble staged at sources/stacks-spaces-cohomology/preamble.tex is excerpted verbatim as three spans, lines 8-12 (verbatim package and the project's own reference/slogan/history display environments), lines 14-18 (Xy-pic setup) and lines 154-217 (theorem declarations and mathematical macros). The wrapper class and the amsmath/amssymb/amsthm/graphicx base packages are supplied by the assembler. No macro was invented or redefined.\par\medskip
\noindent\textit{Editorial note.} Same-chapter pointers are located at their \textbackslash{}label lines inside the pinned chapter file; cross-chapter pointers are located in the project's own tag list staged in the snapshot (sources/stacks-chow/tags.txt), at the verbatim lines carrying their canonical tags.\par\medskip
\noindent\textit{Allowed standard prerequisites (editorial).}\begin{itemize}
\item Flatness of the base changed morphism, and flatness of the induced morphisms of small etale sites. Morphisms of Spaces, labels spaces-morphisms-lemma-base-change-flat (invoked at proof line 1950) and spaces-morphisms-lemma-flat-morphism-sites (invoked at proof line 1953). They show that \$g'\$ is flat and that flatness of \$g\$ and \$g'\$ is the flatness of the corresponding morphisms of small etale ringed sites, which is what the base change map requires.
\item The base change map for a flat morphism of sites. Cohomology on Sites, label sites-cohomology-lemma-base-change-map-flat-case (invoked at proof line 1956). It supplies the map \$g\textasciicircum{}*R\textasciicircum{}pf\_*\textbackslash{}mathcal\{F\} \textbackslash{}to R\textasciicircum{}pf'\_*\textbackslash{}mathcal\{F\}'\$ whose bijectivity is the content of part (1).
\item Vanishing of the higher direct images of the affine morphism \$g'\$, with Leray. Same chapter spaces-cohomology.tex, label lemma-affine-vanishing-higher-direct-images (invoked at proof line 1977) and Cohomology on Sites, label sites-cohomology-lemma-apply-Leray (invoked at proof line 1979). They give \$H\textasciicircum{}p(X\_B, \textbackslash{}mathcal\{F\}\_B) = H\textasciicircum{}p(X, g'\_*\textbackslash{}mathcal\{F\}\_B)\$, because \$g'\$ is affine as a base change of the affine morphism \$\textbackslash{}Spec(B) \textbackslash{}to \textbackslash{}Spec(A)\$.
\item Lazard's theorem, writing a flat algebra as a filtered colimit of finite free modules. Algebra, label algebra-theorem-lazard (invoked at proof line 1997). It provides the filtered system of finite free \$A\$-modules \$M\_i\$ with \$B = \textbackslash{}colim M\_i\$, over which the tensor product and the cohomology can be moved.
\item Cohomology of a quasi-compact quasi-separated space commutes with filtered colimits. Same chapter spaces-cohomology.tex, label lemma-colimits (invoked at proof line 2014). It is the third equality in the chain, and it is what allows the colimit to be moved out of the cohomology of the tensor product.
\end{itemize}\medskip
\ifdefined\hypertarget\hypertarget{sourcepointer-1}{}\fi
\noindent\textit{Source reference.} Original theorem 058G: sources/stacks-chow/tags.txt, source lines 6779--6779. Navigation pointer for \textbackslash{}ref\{algebra-theorem-lazard\} used inside the counted span (line 1997); the argument that consumes it is retained in the span.\par\medskip
\ifdefined\hypertarget\hypertarget{sourcepointer-2}{}\fi
\noindent\textit{Source reference.} Original lemma 073H: sources/stacks-spaces-cohomology/pinned-original-source.tex, source lines 1522--1522. Navigation pointer for \textbackslash{}ref\{lemma-affine-vanishing-higher-direct-images\} used inside the counted span (line 1977); the argument that consumes it is retained in the span.\par\medskip
\ifdefined\hypertarget\hypertarget{sourcepointer-3}{}\fi
\noindent\textit{Source reference.} Original lemma 073E: sources/stacks-spaces-cohomology/pinned-original-source.tex, source lines 667--667. Navigation pointer for \textbackslash{}ref\{lemma-colimits\} used inside the counted span (line 2014); the argument that consumes it is retained in the span.\par\medskip
\ifdefined\hypertarget\hypertarget{sourcepointer-4}{}\fi
\noindent\textit{Source reference.} Original lemma 0733: sources/stacks-chow/tags.txt, source lines 9047--9047. Navigation pointer for \textbackslash{}ref\{sites-cohomology-lemma-apply-Leray\} used inside the counted span (line 1979); the argument that consumes it is retained in the span.\par\medskip
\ifdefined\hypertarget\hypertarget{sourcepointer-5}{}\fi
\noindent\textit{Source reference.} Original lemma 0736: sources/stacks-chow/tags.txt, source lines 9050--9050. Navigation pointer for \textbackslash{}ref\{sites-cohomology-lemma-base-change-map-flat-case\} used inside the counted span (line 1938); the argument that consumes it is retained in the span.\par\medskip
\ifdefined\hypertarget\hypertarget{sourcepointer-6}{}\fi
\noindent\textit{Source reference.} Original lemma 03MO: sources/stacks-chow/tags.txt, source lines 4756--4756. Navigation pointer for \textbackslash{}ref\{spaces-morphisms-lemma-base-change-flat\} used inside the counted span (line 1950); the argument that consumes it is retained in the span.\par\medskip
\ifdefined\hypertarget\hypertarget{sourcepointer-7}{}\fi
\noindent\textit{Source reference.} Original lemma 073C: sources/stacks-chow/tags.txt, source lines 9056--9056. Navigation pointer for \textbackslash{}ref\{spaces-morphisms-lemma-flat-morphism-sites\} used inside the counted span (line 1953); the argument that consumes it is retained in the span.\par\medskip
\begin{lemma}[Flat base change]
\label{lemma-flat-base-change-cohomology}
Let $S$ be a scheme. Consider a cartesian diagram of algebraic spaces
$$
\xymatrix{
X' \ar[d]_{f'} \ar[r]_{g'} & X \ar[d]^f \\
Y' \ar[r]^g & Y
}
$$
over $S$.
Let $\mathcal{F}$ be a quasi-coherent $\mathcal{O}_X$-module
with pullback $\mathcal{F}' = (g')^*\mathcal{F}$.
Assume that $g$ is flat and that $f$ is quasi-compact and quasi-separated.
For any $i \geq 0$
\begin{enumerate}
\item the base change map of
Cohomology on Sites, Lemma
\ref{sites-cohomology-lemma-base-change-map-flat-case}
is an isomorphism
$$
g^*R^if_*\mathcal{F} \longrightarrow R^if'_*\mathcal{F}',
$$
\item if $Y = \Spec(A)$ and $Y' = \Spec(B)$, then
$H^i(X, \mathcal{F}) \otimes_A B = H^i(X', \mathcal{F}')$.
\end{enumerate}
\end{lemma}

\begin{proof}
The morphism $g'$ is flat by
Morphisms of Spaces, Lemma \ref{spaces-morphisms-lemma-base-change-flat}.
Note that flatness of $g$ and $g'$ is equivalent to flatness
of the morphisms of small \'etale ringed sites, see
Morphisms of Spaces, Lemma \ref{spaces-morphisms-lemma-flat-morphism-sites}.
Hence we can apply
Cohomology on Sites, Lemma
\ref{sites-cohomology-lemma-base-change-map-flat-case}
to obtain a base change map
$$
g^*R^pf_*\mathcal{F} \longrightarrow R^pf'_*\mathcal{F}'
$$
To prove this map is an isomorphism we can work locally in the \'etale
topology on $Y'$. Thus we may assume that $Y$ and $Y'$ are affine
schemes. Say $Y = \Spec(A)$ and $Y' = \Spec(B)$.
In this case we are really trying to show that the map
$$
H^p(X, \mathcal{F}) \otimes_A B \longrightarrow H^p(X_B, \mathcal{F}_B)
$$
is an isomorphism where $X_B = \Spec(B) \times_{\Spec(A)} X$ and
$\mathcal{F}_B$ is the pullback of $\mathcal{F}$ to $X_B$.
In other words, it suffices to prove (2).

\medskip\noindent
Fix $A \to B$ a flat ring map and let $X$ be a quasi-compact and
quasi-separated algebraic space over $A$. Note that $g' : X_B \to X$
is affine as a base change of $\Spec(B) \to \Spec(A)$. Hence
the higher direct images $R^i(g')_*\mathcal{F}_B$ are zero by
Lemma \ref{lemma-affine-vanishing-higher-direct-images}.
Thus $H^p(X_B, \mathcal{F}_B) = H^p(X, g'_*\mathcal{F}_B)$, see
Cohomology on Sites, Lemma \ref{sites-cohomology-lemma-apply-Leray}.
Moreover, we have
$$
g'_*\mathcal{F}_B = \mathcal{F} \otimes_{\underline{A}} \underline{B}
$$
where $\underline{A}$, $\underline{B}$ denotes the constant sheaf of
rings with value $A$, $B$. Namely, it is clear that there is a map
from right to left. For any affine scheme $U$ \'etale over $X$ we have
\begin{align*}
g'_*\mathcal{F}_B(U) & = \mathcal{F}_B(\Spec(B) \times_{\Spec(A)} U) \\
& =
\Gamma(\Spec(B) \times_{\Spec(A)} U,
(\Spec(B) \times_{\Spec(A)} U \to U)^*\mathcal{F}|_U) \\
& =
B \otimes_A \mathcal{F}(U)
\end{align*}
hence the map is an isomorphism. Write $B = \colim M_i$ as a filtered
colimit of finite free $A$-modules $M_i$ using Lazard's theorem, see
Algebra, Theorem \ref{algebra-theorem-lazard}.
We deduce that
\begin{align*}
H^p(X, g'_*\mathcal{F}_B) &
= H^p(X, \mathcal{F} \otimes_{\underline{A}} \underline{B}) \\
& = H^p(X, \colim_i \mathcal{F} \otimes_{\underline{A}} \underline{M_i}) \\
& = \colim_i H^p(X, \mathcal{F} \otimes_{\underline{A}} \underline{M_i}) \\
& = \colim_i H^p(X, \mathcal{F}) \otimes_A M_i \\
& = H^p(X, \mathcal{F}) \otimes_A \colim_i M_i \\
& = H^p(X, \mathcal{F}) \otimes_A B
\end{align*}
The first equality because
$g'_*\mathcal{F}_B = \mathcal{F} \otimes_{\underline{A}} \underline{B}$
as seen above.
The second because $\otimes$ commutes with colimits.
The third equality because cohomology on $X$ commutes with
colimits (see
Lemma \ref{lemma-colimits}).
The fourth equality because $M_i$ is finite free (i.e., because cohomology
commutes with finite direct sums).
The fifth because $\otimes$ commutes with colimits.
The sixth by choice of our system.
\end{proof}

\end{document}
